% ============================================================
%  Archivo: seccion_problema.tex
%  Uso: pegar en tu proyecto Overleaf o subir como archivo nuevo
%  Compilar con: pdflatex + bibtex (o biblatex)
% ============================================================

\documentclass[12pt,a4paper]{article}

% ── Paquetes ─────────────────────────────────────────────────
\usepackage[utf8]{inputenc}
\usepackage[T1]{fontenc}
\usepackage[spanish]{babel}
\usepackage{geometry}
\geometry{a4paper, top=2.5cm, bottom=2.5cm, left=3cm, right=2.5cm}

\usepackage{setspace}
\onehalfspacing                  % interlineado 1.5

\usepackage{parskip}             % espacio entre párrafos
\setlength{\parindent}{0pt}

\usepackage{hyperref}
\hypersetup{colorlinks=true, linkcolor=black, citecolor=black, urlcolor=blue}

\usepackage[backend=bibtex, style=ieee, sorting=none]{biblatex}
\addbibresource{referencias.bib}  % <- nombre de tu archivo .bib

% ── Metadatos ────────────────────────────────────────────────
\title{\textbf{Problema de Investigación} \\
       \large Sistema de Gestión Nutricional basado en RAG}
\author{Kevin Patricio Sarango Olaya \\
        \small Proyectos Tecnológicos 2 --- Ciclo 8 ``A'' \\
        \small FEIRNNR --- Universidad Nacional de Loja}
\date{Mayo 2026}

% ============================================================
\begin{document}
\maketitle
\clearpage

% ============================================================
\section{Problema de investigación}
% ============================================================

El problema de investigación se define como la descripción clara, objetiva y
contextualizada de una realidad que genera una necesidad, conflicto o desafío
en un entorno determinado. En este sentido, el problema debe sustentarse en
hechos concretos, tales como datos, observaciones o evidencias provenientes de
fuentes confiables, evitando juicios de valor o generalizaciones.

Asimismo, es fundamental ubicar el problema dentro de un contexto específico,
identificando dónde ocurre y a quiénes afecta, lo que permite delimitar
adecuadamente el alcance de la investigación. En el ámbito de la salud
nutricional, esto implica considerar tanto a los profesionales de la salud como
a los pacientes que requieren seguimiento continuo y personalizado.

De esta manera, la formulación del problema de investigación constituye la base
para el desarrollo del estudio, ya que orienta la definición de objetivos,
hipótesis y metodología, garantizando la coherencia y pertinencia del proyecto.

% ------------------------------------------------------------
\subsection{Situación Problemática}
% ------------------------------------------------------------

En Ecuador, el seguimiento de la salud nutricional enfrenta barreras como la
falta de tiempo y el desconocimiento del estado físico real del paciente, lo
que evidencia la necesidad de herramientas tecnológicas que complementen la
consulta profesional~\cite{vargas2025}. La gestión nutricional en centros
especializados sigue siendo predominantemente manual, utilizando hojas de
cálculo o registros en papel, lo que genera vulnerabilidad en los datos,
ralentiza la atención y dificulta el acceso rápido a la información histórica
de los pacientes. Esta falta de digitalización incrementa los errores en la
planificación y obstaculiza el seguimiento de la evolución fisiológica del
paciente, debilitando el vínculo médico-paciente~\cite{quichimbo2021}.
A día de hoy, las herramientas digitales disponibles carecen de interfaces
intuitivas adaptadas al contexto ecuatoriano, lo que provoca el abandono del
seguimiento ante planes estáticos y una retroalimentación profesional
insuficiente~\cite{vargas2025,quichimbo2021,siddiqui2025}.

La gestión digital de planes nutricionales se fundamenta en los principios de
la Nutrición Personalizada~(NP), la cual propone utilizar información
específica del paciente (patológica, fisiológica y hereditaria) para diseñar
intervenciones dietéticas más efectivas que las guías
generales~\cite{siddiqui2025}. Este enfoque busca superar las limitaciones de
las recomendaciones tradicionales mediante el uso de tecnologías que permitan
analizar datos individuales de manera eficiente. Los Sistemas de Soporte a la
Decisión Clínica permiten integrar y procesar datos nutricionales de forma
ágil, facilitando la toma de decisiones del profesional de salud y reduciendo
errores en la planificación dietética~\cite{quichimbo2021}. Asimismo, las
plataformas de salud móvil~(mHealth) fortalecen la interacción entre paciente
y especialista, promoviendo un seguimiento continuo y una mayor adherencia a
los planes nutricionales~\cite{ghose2021}. La personalización efectiva de los
sistemas nutricionales requiere el uso de bases de datos adaptadas al entorno
local, garantizando pertinencia de las recomendaciones alimentarias en
contextos específicos~\cite{aldarras2025}.

En los artículos~\cite{cunha2022} y~\cite{theodore2024}, los autores
demostraron que dispositivos IoT reducen la carga manual de registro en
pacientes, aunque persisten limitaciones relacionadas con la precisión y la
privacidad. Las investigaciones~\cite{lee2023} y~\cite{mortazavi2023}
validaron que plataformas basadas en inteligencia artificial junto con
monitores continuos de glucosa reportaron reducciones en la~HbA1c
(hemoglobina glicosilada), aunque evidenciaron errores del~22\% en la
predicción de macronutrientes. En los artículos~\cite{gavai2025}
y~\cite{arshad2025}, los autores implementaron sistemas de generación
aumentada por recuperación~(RAG) con modelos de lenguaje locales, alcanzando
hasta un~80\% de adherencia y ajustes dinámicos, aunque con una validación
clínica aún incipiente. Los autores de los artículos~\cite{singar2024}
y~\cite{xu2025} describen variantes genéticas~(FTO, APOE) y perfiles de
microbiota como moduladores de la respuesta dietética, aunque su integración
en plataformas digitales es parcial. En los artículos~\cite{abeltino2025}
y~\cite{ebadi2025}, los autores señalan que la falta de interoperabilidad, la
escasa adaptación regional y los altos costos limitan la escalabilidad de
estas soluciones.

El estudio se enfoca en la integración de sistemas inteligentes en el
desarrollo de software, permitiendo una transición hacia un modelo proactivo
en la gestión de la salud nutricional, alineándose con el ODS~9 (Industria,
Innovación e Infraestructura), el cual promueve el fortalecimiento de la
infraestructura tecnológica y la innovación como pilares para mejorar la
eficiencia de los servicios de salud~\cite{un2024}. La línea de investigación
en la que se enmarca este estudio está orientada al desarrollo tecnológico, la
innovación y la generación de soluciones basadas en software para problemáticas
reales, contribuyendo a la transformación digital del sector mediante el uso de
herramientas avanzadas de análisis de datos~\cite{unl2021}. La incorporación de
un modelo de lenguaje permitirá asistir al nutricionista en el procesamiento de
información del paciente, generando planes personalizados y dinámicos,
reduciendo la carga operativa y mejorando la precisión en la toma de
decisiones, fortaleciendo así la adherencia del paciente mediante un
seguimiento continuo.

El propósito del autor es contribuir al fortalecimiento de la gestión
nutricional mediante el desarrollo de una herramienta tecnológica que apoye al
nutricionista como núcleo central del sistema. La propuesta se orienta a
mejorar las funcionalidades de los sistemas existentes, incorporando
características adaptadas al contexto nacional, como el registro de alimentos
locales y la gestión eficiente de la información del paciente, apoyándose en
PostgreSQL como sistema gestor de base de datos relacional, el cual garantiza
la integridad, persistencia y consulta eficiente de los datos clínicos y
nutricionales almacenados. Asimismo, se busca innovar mediante la integración
de un modelo de lenguaje como herramienta de apoyo, que facilite el
procesamiento de datos y la generación de recomendaciones personalizadas a
partir de la información estructurada en la base de datos; con ello, se
pretende reducir las limitaciones presentes en sistemas actuales y mejorar la
calidad del servicio nutricional, fortaleciendo la toma de decisiones clínicas.

% ------------------------------------------------------------
\subsection{Pregunta de investigación}
% ------------------------------------------------------------

¿Qué nivel de satisfacción global reporta el nutricionista con las
recomendaciones de planes personalizados generadas por un prototipo de sistema
basado en RAG, utilizando datos antropométricos y bioquímicos del paciente
junto con un catálogo alimentario ecuatoriano?

% ------------------------------------------------------------
\subsection{Justificación}
% ------------------------------------------------------------

En cuanto al aporte al conocimiento, esta investigación contribuye al campo de
la ingeniería de software y la informática en salud mediante el diseño de un
prototipo que integra un modelo largo de lenguaje con técnicas de generación
aumentada por recuperación~(RAG), posicionando al nutricionista como el actor
central del proceso clínico. En el contexto ecuatoriano, las propuestas
actuales de gestión nutricional carecen de herramientas inteligentes que
apoyen al profesional en la selección de recomendaciones dietéticas
personalizadas a partir de datos del
paciente~\cite{vargas2025,quichimbo2021}. Por ello, el desarrollo de este
prototipo constituirá un aporte técnico concreto al campo de los sistemas de
soporte en la decisión clínica en el contexto ecuatoriano, integrando
procesamiento de datos antropométricos y bioquímicos con un catálogo de
alimentos locales.

Desde una perspectiva científica, el estudio se fundamenta en la transición de
un modelo de gestión nutricional reactivo hacia uno proactivo, sustentado en la
Nutrición Personalizada~(NP) y la medicina de
precisión~\cite{siddiqui2025,singar2024}. El sistema sugiere al nutricionista
un conjunto de recomendaciones fundamentadas en bases de conocimientos
validadas y bases de datos nutricionales estructuradas, preservando en todo
momento la confidencialidad de la información del
paciente~\cite{gavai2025,arshad2025}. La decisión final sobre el plan
nutricional recae siempre en el profesional, lo que aporta rigor clínico al
proceso y contribuye a solucionar las limitaciones propias de la planificación
manual, que es propensa a errores humanos y dificulta el seguimiento continuo
de la evolución fisiológica del paciente~\cite{mortazavi2023}.

En el ámbito social, el proyecto responde a la necesidad de mejorar el
bienestar nutricional de la población ecuatoriana, beneficiando directamente a
los nutricionistas de centros de salud y consultorios privados, quienes
contarían con una herramienta de apoyo para la selección informada de
recomendaciones dietéticas; de forma indirecta, los pacientes se beneficiarán
al recibir planes de mayor calidad y pertinencia, validados por el profesional
con base en sugerencias del sistema fundamentadas en datos del
paciente~\cite{ghose2021}. Al fortalecer la capacidad de decisión del
nutricionista mediante una plataforma digital, se busca incrementar la calidad
de los planes nutricionales propuestos y evaluar su factibilidad de
implementación, contribuyendo así a mejorar la atención clínica de los
pacientes~\cite{abeltino2025,ebadi2025}.

Finalmente, en el aspecto económico, la investigación se vincula con el~ODS~9
al fortalecer la infraestructura tecnológica e impulsar la innovación en los
servicios de salud~\cite{un2024}. La implementación de esta herramienta
reducirá la carga operativa del profesional, permitiéndole centrarse en la
evaluación clínica y en la selección informada del plan más
adecuado~\cite{quichimbo2021}. Al trabajar únicamente con los datos bioquímicos
y antropométricos del paciente, la solución reduce los riesgos legales y éticos
asociados al manejo de información sensible en salud, promoviendo una
infraestructura digital segura y escalable. Este estudio se enmarca en las
líneas de investigación aprobadas de la Universidad Nacional de Loja,
orientadas al desarrollo tecnológico, la innovación y la generación de
soluciones basadas en software para problemáticas reales del entorno nacional,
contribuyendo a la transformación digital del sector salud~\cite{unl2021}.

% ============================================================
%  BIBLIOGRAFÍA
% ============================================================
\clearpage
\printbibliography[title={Referencias}]

\end{document}
